\subsection{Mais Indução}

\begin{enumerate}[label=2.\arabic*.]
    
    \item (Pequeno Teorema de Fermat) Se $p$ é um primo qualquer, mostre usando a fórmula binomial 
    que $n^p - n$ é multíplo de $p$ para cada inteiro $n$.  
    
    \item (Média aritmética e média geométrica) Dados $a_1, a_2, \ldots, a_n$ inteiros $n > 0$, se 
    $A = \frac{a_1 + a_2 + \ldots + a_n}{n}$ e $G = \sqrt[n]{a_1a_2 \ldots a_n}$, prove por indução
    que $G \le A$ \\
    (Sugestão: supondo $a_1 \le a_2 \le \ldots \le a_n$, use $a_1 + a_n \ge \frac{a_1a_n}{G} + G$)
    
    \item (Levantando pesos). Dados uma balança de dois pratos e $n$ pesos de
    $1, 3, \ldots, 3^{n-1}\,\text{kg}$, prove que, ao colocar-se pesos nos
    pratos, é possível pesar $N$ kg, onde $N$ é um inteiro $\le 1$ e $\le \frac{3^n - 1}{2}$.
    (Dica: considere as somas da forma
    $3^{n-1}e_{n-1} + \cdots + 3e_1 + e_0$, com $e_i \in \{-1, 0, +1\}$.)

    \item (Polinômios em várias indeterminadas). 
    Demonstre que o número de termos em cada polinômio homogêneo de grau $d$
    em $n$ indeterminadas é, no máximo, $\frac{(n + d)!}{n!d!}$
    (Dica: indução sobre $n$; observe que o número de termos em um polinômio
    homogêneo de grau $d$ em $n$ indeterminadas é o mesmo que de um polinômio
    de grau $d$ em $n-1$ indeterminadas.)
    
    \item (Um famoso determinante). 
    Calcule o determinante da matriz
    \[
    \begin{pmatrix}
    1 & x_1 & x_1^2 & \cdots & x_1^{n-1} \\
    1 & x_2 & x_2^2 & \cdots & x_2^{n-1} \\
    \vdots & \vdots & \vdots & \ddots & \vdots \\
    1 & x_n & x_n^2 & \cdots & x_n^{n-1}
    \end{pmatrix}.
    \]
    (Dica: na última linha, considerando $x_n$ como uma indeterminada,
    obtenha um polinômio de grau $n-1$ que se anula em $x_1, \ldots, x_{n-1}$.)
    
    \item (Grafos I).
    Um \emph{grafo} consiste num conjunto $\Gamma$ com um subconjunto distinguido
    não vazio $V$ e duas aplicações $d_0, d_1 : \Gamma \to V$ tais que
    $d_0(v) = v = d_1(v)$ para cada $v \in V$. Os vértices de $\Gamma$ são os elementos 
    de $V$; as arestas são os elementos de $E = \Gamma - V$.
    Para cada aresta $e$ de $\Gamma$ introduz-se símbolos formais $e^{\pm 1}$, de modo que
    $d_0(e^1) = d_0(e) = d_1(e^{-1})$ e $d_1(e^1) = d_1(e) = d_0(e^{-1})$. 
    Um \emph{caminho} entre vértices $u$ e $w$ de $\Gamma$ é uma sequência finita
    $e_1^{\epsilon_1}, \ldots, e_m{\epsilon_m}$, em que $m \ge 0$, e $e_i \in E$,
    $\epsilon_i = \pm 1$, $d_0(e_1^{\epsilon_1}) = u$, $d_1(e_m^{\epsilon_m}) = w$
    e $d_1(e_i^{\epsilon_i}) = d_0(e_{i+1}^{\epsilon_{i + 1}})$;
    tal caminho é reduzido se $e_i^{\epsilon_i} \neq e_{i+1}^{-\epsilon_{i+1}}$ 
    quando $1 \le i < m$. Se para cada $v \in V$ existe um único caminho reduzido entre $v$ e
    $v$, e ainda, para qualquer $w \in V$, existe ao menos um caminho entre
    $w$ e $v$, diz-se que $\Gamma$ é uma árvore.
    Prove que, se $\Gamma$ é finito e uma árvore, então a quantidade de seus
    vértices menos a quantidade de suas arestas é $1$.
\end{enumerate}
